% Lösungen Einheit 3 von 5 – Prozentwert berechnen (zusammenbau v0.1)
\einheitenkopf{Einheit 3 von 5 · Prozentwert berechnen}

%% TODO Lösung der Erklärzeile 17f) fehlt in der Bank
\erg{17}{a) $40\,€$ \quad b) $15$ kg \quad c) $30$ m \quad d) $7$ l \quad e) $9\,€$ \quad f) \ldots \quad g) $1\,\% = 4$ kg; $6\,\% = 24$ kg \quad h) $10\,\% = 9\,€$; $30\,\% = 27\,€$ \quad i) $0{,}45 \cdot 80 = 36$ kg \quad j) $0{,}2 \cdot 4{,}50 = 0{,}90\,€$ \quad k) Ersparnis $0{,}2 \cdot 70 = 14\,€$ \quad l) $1{,}5 \cdot 40 = 60$ kg \quad m) $72\,€$ ($480 \cdot 0{,}15 = 72$; $408\,€$ ist der Restpreis) \quad n) $60 \cdot 0{,}17 = 10{,}20\,€$ \quad o) $52\,\% - 35\,\% = 17$ Prozentpunkte; $0{,}17 \cdot 800 = 136$ Haushalte (nicht $17$)}
\erg{18}{a) $0{,}19 \cdot 300 = 57\,€$}
\erg{19}{a) Fehler: $60\,\%$ als $0{,}06$ geschrieben. Richtig: $0{,}6 \cdot 35 = 21\,€$. \quad b) Das Ganze sind $100\,\%$; in $100$ gleiche Teile geteilt ist jeder Teil $1\,\%$, also $500 : 100$. \quad c) $18\,€$ ($3$ von $10$ Teilen) \quad d) A: $95 - 19 = 76\,€$; B: $85 - 8{,}50 = 76{,}50\,€$ – in Laden A, um $0{,}50\,€$.}
